% from ega2-1-fr.tex
\oldpage[II]{134}

de $X$. Nous allons voir que si $t$ est la section $s$ considérée comme
section de $\mathcal{M}^{\otimes n}$ au-dessus de $Z$, $t$ répond à la
question. Posons en effet
\[
\label{II.6.6.9.2-fr}
V=X-g(X'-X'_{s'})
\tag{6.6.9.2}
\]
qui est un ouvert de $X$ contenant $x$ et contenu dans $U$, en vertu de
(\hyperref[II.6.6.9.1-fr]{6.6.9.1}) et
(\hyperref[II.6.1.10-fr]{6.1.10}). Nous allons montrer que l'on a
\[
\label{II.6.6.9.3-fr}
h^{-1}(V)\subset Z_t\subset h^{-1}(U),
\tag{6.6.9.3}
\]
ce qui achèvera la démonstration. Il revient au même de dire que l'ensemble
$T$ des $y\in X$ tels que $s_y$ soit inversible, contient $V$ et est contenu
dans $U$. Pour cela, considérons d'abord un ouvert affine $W$ contenu dans
$V$~; $g^{-1}(W)$ est ouvert affine dans $X'$ et en vertu de
(\hyperref[II.6.6.9.2-fr]{6.6.9.2}), $s'_{y'}$ est inversible pour tout
$y'\in g^{-1}(W)$~; en vertu des hypothèses sur $X$ et $\mathcal{B}$, on
peut appliquer les résultats de (\hyperref[II.6.4.7-fr]{6.4.7}) et on voit
que si $y=g(y')$, $s_y$ est inversible, autrement dit $V\subset T$. D'autre
part, il résulte aussi de (\hyperref[II.6.4.7-fr]{6.4.7}) que, inversement,
si $s_y$ est inversible, il en est de même de $s'_{y'}$, ce qui implique
$y'\in g^{-1}(U)$ par (\hyperref[II.6.6.9.1-fr]{6.6.9.1}) et par suite
$y\in U$, d'où $T\subset U$ dans ce cas.

On passe de là au cas général par le même raisonnement que dans
(\hyperref[II.6.6.2-fr]{6.6.2}), remplaçant $X'$ par $X''$ tel que
\[
g'_*(\mathcal{O}_{X''})\longrightarrow
g'_*(\mathcal{O}_{X''})\otimes_{\mathcal{O}_X}\mathcal{R}(X)
\]
soit injectif, et $\mathcal{L}'$ par un $\mathcal{O}_{X''}$-Module ample
$\mathcal{L}''$ tel que
$N'(g_*(\mathcal{L}'))=N'(g'_*(\mathcal{L}''))$. Le lemme
(\hyperref[II.6.6.9-fr]{6.6.9}) est donc démontré dans tous les cas (avec un
choix convenable de $h$).

\begin{corollary}[6.6.10]
\label{II.6.6.10-fr}
Supposons vérifiées les hypothèses de
(\hyperref[II.6.6.9-fr]{6.6.9})~; pour tout $\mathcal{O}_X$-Module inversible
$\mathcal{L}$ tel que $g^*(\mathcal{L})$ soit ample, $h^*(\mathcal{L})$ est
ample.
\end{corollary}

En effet, si on pose $\mathcal{L}'=g^*(\mathcal{L})$, on a alors
$g_*(\mathcal{L}')=\mathcal{L}\otimes_{\mathcal{O}_X}\mathcal{B}$
(\hyperref[0.5.4.10-fr]{0, 5.4.10}), donc
\[
N'(g_*(\mathcal{L}'))=
(\mathcal{L}\otimes_{\mathcal{O}_X}\mathcal{C})^{\otimes n}
\]
(par le même raisonnement que pour
\hyperref[II.6.5.2.4-fr]{6.5.2.4}). On en conclut que l'on a
$\mathcal{M}=(h^*(\mathcal{L}))^{\otimes n}$, et comme $\mathcal{M}$ est
ample, il en est de même de $h^*(\mathcal{L})$
(\hyperref[II.4.5.6-fr]{4.5.6}).

\begin{proposition}[6.6.11]
\label{II.6.6.11-fr}
Soient $Y$ un schéma affine, $X$ un préschéma noethérien réduit,
$f:X\to Y$ un morphisme quasi-compact, $g:X'\to X$ un morphisme fini
surjectif. Soit $W$ une partie ouverte de $X$ telle que, pour tout $x\in W$,
ou bien $X$ est normal au point $x$, ou bien il existe un voisinage ouvert
$T\subset W$ de $x$ tel que $(g_*(\mathcal{O}_{X'}))|T$ soit un
$(\mathcal{O}_X|T)$-Module localement libre. Il existe alors un
$Y$-préschéma réduit $Z$ et un $Y$-morphisme fini et birationnel
$h:Z\to X$ tel que la restriction de $h$ à $h^{-1}(W)$ soit un isomorphisme
$h^{-1}(W)\xrightarrow{\sim}W$ et qui possède la propriété suivante~: pour
tout $\mathcal{O}_X$-Module inversible $\mathcal{L}$ tel que
$g^*(\mathcal{L})$ soit ample relativement à $f\circ g$, $h^*(\mathcal{L})$
est ample relativement à $f\circ h$.
\end{proposition}

Comme $Y$ est affine, $g^*(\mathcal{L})$ est ample, et il s'agit alors de
prouver que pour un choix convenable de $h$, $h^*(\mathcal{L})$ est ample
(\hyperref[II.4.6.6-fr]{4.6.6}). Nous allons voir qu'on peut remplacer $g$
par un morphisme fini surjectif $g':X''\to X$ tel que
${g'}^*(\mathcal{L})$ soit ample et que
$g'_*(\mathcal{O}_{X''})\otimes_{\mathcal{O}_X}\mathcal{R}(X)$ soit un
$\mathcal{R}(X)$-Module \emph{localement libre de rang constant}~; on sera
alors ramené aux conditions de (\hyperref[II.6.6.10-fr]{6.6.10}) et la
proposition sera démontrée.

Pour cela, soit $\mathcal{B}=g_*(\mathcal{O}_{X'})$~; désignons par $X_i$
($1\leq i\leq n$) les sous-préschémas fermés réduits de $X$ qui ont pour
espaces sous-jacents les composantes irréductibles de $X$
(\hyperref[I.5.2.1-fr]{I, 5.2.1})~; ils sont intègres par hypothèse. Soient
$X'_i$ le sous-préschéma fermé

